\documentclass[11pt]{article}
\usepackage[utf8]{inputenc}
\usepackage[T1]{fontenc}
\usepackage{lmodern}
\usepackage{amsmath,amssymb}
\usepackage{microtype}
\usepackage[letterpaper,margin=1.1in]{geometry}
\usepackage{booktabs}
\usepackage[hidelinks]{hyperref}
\usepackage{cleveref}

\newcommand{\fhat}{\hat{f}}
\newcommand{\enc}{\operatorname{enc}}
\newcommand{\dec}{\operatorname{dec}}
\newcommand{\B}{\{0,1\}}

\title{Trapdoor Computing:\\[0.2em]\large Computing on Opaque Encodings}
\author{Alexander Towell\\\texttt{lex@metafunctor.com}}
\date{2026}

\begin{document}
\maketitle

\begin{abstract}
Imagine a service that cannot read your queries and cannot read its own
answers, yet returns the right ones. We make that precise. \emph{Trapdoor
computing} is a paradigm in which a \emph{trusted} machine holds an encoder
and a decoder while an \emph{untrusted} machine sees only total functions on
opaque bit strings flowing through opaque lookup tables. The untrusted machine
cannot tell a real query from filler, one encoding from another, or a correct
result from noise; its privacy comes from a one-way hash and a uniform
representation, not from hiding access patterns. This is deliberately not
oblivious RAM, not fully homomorphic encryption, and not a simulation-based
security game. Confidentiality here is \emph{measured} in the sense of
quantitative information flow, an explicit leakage quantity a designer can
compute and trade against space, bandwidth, and accuracy, rather than assumed
\emph{negligible} in the sense of a cryptographic reduction. This paper states
the paradigm, draws its boundary against the cryptographic tools it is not,
and maps the body of work that develops it: an abstraction, an algebra, a
confidentiality theory, and concrete constructions.
\end{abstract}

\section{The Oblivious Service}
\label{sec:service}

Consider a server that answers membership queries against a private set: a
watchlist, a revocation list, an index. We want it to answer correctly while
learning nothing it could turn against us, not which name was asked, not which
names are on the list, not even how often a given name is asked. The standard
cryptographic tools can do versions of this at a price: oblivious RAM hides the
access pattern at a polylogarithmic bandwidth cost per access; fully
homomorphic encryption computes on ciphertexts at a large constant and circuit-%
depth cost; searchable symmetric encryption is fast but leaks a documented
profile that inference attacks exploit. Each makes a different trade, and each
is the right tool for some regime.

Trapdoor computing occupies a different point. A \emph{trusted} machine $T$
holds a secret and the operations to encode an input and decode an output. An
\emph{untrusted} machine $U$ holds a single object, a total function $\fhat$ on
bit strings, and evaluates it on whatever it is given. To answer a query, $T$
encodes the input as a token $c$, $U$ returns $\fhat(c)$, and $T$ decodes the
result. The untrusted machine never holds the secret, never decodes, and never
sees a value, only opaque tokens flowing through an opaque table. The privacy
is not access-pattern hiding; it is that the tokens carry no readable signal.

\section{The Paradigm}
\label{sec:paradigm}

The organizing object is the \emph{cipher map}: a total function $\fhat : \B^n
\to \B^n$ that implements a trapdoor approximation of a latent function
$f : X \to Y$, paired with an encoder $\enc$, a decoder $\dec$, and a secret
$s$ from which all three derive~\cite{towell2026cipher}. The trusted machine
computes $\dec(\fhat(\enc(x)))$ and reads $f(x)$; the untrusted machine
computes $\fhat$ on tokens and reads nothing. Four features of the definition
do the work.

\paragraph{Totality.} $\fhat$ is defined on \emph{all} of $\B^n$, not only on
the tokens that encode real inputs. Hand it a string that encodes nothing and
it still returns an $n$-bit answer. There is no ``not found,'' no exception, no
out-of-band signal. This is the first privacy guarantee: with no special case
to observe, the untrusted machine cannot separate meaningful queries from
filler.

\paragraph{Representation uniformity.} A value may be encoded as any of several
tokens, and the encoding is chosen so that the stream of tokens the untrusted
machine sees is close, in total variation, to uniform on the populated support.
Frequent values are spread across more representations, classical homophonic
substitution, so that no token is queried often and the traffic betrays nothing
about which value was asked.

\paragraph{Correctness.} The map is correct up to an error budget $\eta$: it
computes $f$ on in-domain inputs with probability at least $1 - \eta$. A nonzero
$\eta$ is not only a construction cost. It is also privacy: a reported result
may be a true answer or a cipher-map error, and the ambiguity is deniability
the untrusted machine cannot resolve.

\paragraph{Composability.} Cipher maps chain. The output token of one is a well-%
formed input to the next, and the error of a chain is bounded predictably by
the errors of its links. This is what makes trapdoor computing a calculus
rather than a collection of isolated lookups.

\paragraph{The four cannots.} Against the untrusted machine, these properties
guarantee four things it cannot do, each blocked by a named property: it cannot
\emph{decode} a token to its value (one-way hash); it cannot \emph{distinguish
a real query from filler} (totality); it cannot \emph{determine the domain} or
read which value a token carries (representation uniformity); and it cannot
\emph{be sure a result is correct} (nonzero $\eta$). These are the paradigm's
promises, and the rest of the program makes each one quantitative.

\section{What Trapdoor Computing Is Not}
\label{sec:not}

The paradigm is easy to mistake for a neighbor, and the mistakes matter because
they import the wrong threat model. Privacy here comes from one-way hashing and
uniform representation, \emph{not} from hiding access patterns; any argument
that rests on access-pattern indistinguishability has drifted into the wrong
framework. \Cref{tab:not} draws the boundary.

\begin{table}[t]
\centering
\caption{Trapdoor computing against its neighbors, by the difference in threat
model, not a claim of superiority.}
\label{tab:not}
\begin{tabular}{@{}l p{3.4cm} p{4.7cm}@{}}
\toprule
Tool & Hides by & Trapdoor computing instead \\
\midrule
Oblivious RAM~\cite{goldreich1996software} & access-pattern indistinguishability & opaque tokens; access pattern not the secret \\
Fully homomorphic enc.~\cite{gentry2009fully} & computing on ciphertext & a hash table, not a circuit; bounded error \\
Garbled circuits & per-evaluation garbling & a reusable total function on tokens \\
SSE / PPE~\cite{curtmola2006searchable,naveed2015inference} & a negligible-or-documented leakage profile & a \emph{measured}, tunable leakage \\
\bottomrule
\end{tabular}
\end{table}

The framework is also not a simulation-based or game-based cryptographic
construction. It makes no claim that an adversary's advantage is negligible in
a security parameter. Its guarantees are leakage quantities, which is the next
point and the heart of the paradigm.

\section{Measured, Not Negligible}
\label{sec:measured}

Cryptography's strongest tools promise that what an adversary learns is
\emph{negligible}: a function vanishing faster than any inverse polynomial in a
security parameter. That promise is the right one when it can be kept at an
affordable cost. Trapdoor computing makes a different and more modest promise,
borrowed from quantitative information flow~\cite{smith2009foundations,%
alvim2020science}: what the adversary learns is a \emph{measured} quantity, a
number a designer can compute and trade.

Concretely, the leakage of a single token is captured by an entropy ratio
$e = H(Q)/H^{*}$, where $H^{*} = \log_2|\mathrm{im}(\enc)|$ is the maximum
entropy on the populated support (the usable bits of a token) and $H(Q)$ the
entropy actually observed; the representation-uniformity parameter $\delta$
lower-bounds it through a Fannes--Audenaert continuity inequality, image-relative
and linear in $\delta$~\cite{towell2026maxconf}. Small
$\delta$ buys proportional confidentiality, at a known cost in space or
bandwidth. This is the sense in which the paradigm's privacy is engineering: not
a guarantee of nothing learned, but an explicit budget of how much, traded
against the resources a system can spend.

The honesty cuts both ways. A measured guarantee is weaker than a negligible
one, and a reviewer schooled in simulation-based definitions will note its
absence. But for the systems that cannot afford oblivious RAM or homomorphic
encryption, a service that cannot read its queries or answers, returns the
right ones, and leaks an amount you can write down is a real and useful
guarantee, and a more honest one than a negligibility claim a construction
cannot actually support.

\section{The Program}
\label{sec:program}

The paradigm has been developed into a body of work, of which this paper is the
front door. The \emph{abstraction} fixes the cipher map and its four measurable
properties, and shows that a single device, the acceptance predicate, unifies
the constructions and traces a space-accuracy duality~\cite{towell2026cipher}.
The \emph{algebra and types} ask which type constructors survive the trapdoor:
products pass, sums hit a genuine impossibility (tag-hiding and untrusted
dispatch are exclusive), and an orbit-closure bound measures what an active
adversary learns by operating on tokens~\cite{towell2026algebraic}. The
\emph{confidentiality theory} measures leakage at two scales that do not reduce
to each other, a marginal scale governed by $\delta$ and a compositional scale,
the joint distribution recoverable from shared-token observations, that no per-%
map parameter can improve~\cite{towell2026maxconf}. And the \emph{constructions}
realize cipher maps concretely, including a linear-algebraic retrieval structure
whose frequency-hiding is structural, holding at zero per-query cost where every
online defense pays per access~\cite{towell2026codec}.

The synthesis of these into one treatment, with the open problems they raise, is
a monograph in preparation. This paper's contribution is the paradigm itself:
the trusted/untrusted model, the cipher map and its four cannots, the boundary
against the cryptographic neighbors, and the stance that confidentiality is
measured rather than negligible.

\section{Conclusion}

Trapdoor computing is a narrow and honest idea: compute on values sealed behind
a one-way trapdoor, let an untrusted machine evaluate opaque tokens through an
opaque table, and measure, rather than assume away, what leaks. It is not the
strongest privacy tool, and it does not try to be. It is the tool for the
middle ground between plaintext and the heavy cryptographic machinery, where a
designer wants a guarantee that is real, composable, and affordable, and is
willing to read the number that says how much.

\paragraph{Provenance.} The framework was begun in 2017 and developed through
2020, then set aside and recovered, re-derived, and carried further in 2024. It
builds on the Bernoulli model of approximate data structures developed
alongside it.

\bibliographystyle{plain}
\bibliography{references}

\end{document}
